\chapter{Analisis}
\label{chap:analisis}

\section{Requisitos}
\label{sec:requisitos}

En este apartado se presentarán los requisitos del sistema desarrollado para este proyecto, definiendo con 
claridad las capacidades y comportamientos que ofrecerá la solución final. Para estructurar su descripción, 
se distinguirán dos categorías fundamentales:

\subsection{Requisitos funcionales}
\label{sec:req_no_funcionales}

Los requisitos funcionales describen las funciones y los comportamientos que el sistema debe ser capaz de realizar 
para cumplir los objetivos especificados. 
Para este proyecto se definieron los siguientes:

\begin{itemize}
    \item \textbf{Captura de movimiento del robot y del objetivo}:
    El sistema deberá ser capaz de obtener la posición del robot y del objetivo de forma correcta. Se usará el sistema de captura Optitrack para la 
    obtención de estos datos.
 
    \item \textbf{Transmisión de datos}:
    El sistema deberá soportar la transmisión de datos desde el Motive al ordenador donde se ejecutan los scripts de control y entrenamiento. La 
    transferencia deberá ser estable, garantizando que las observaciones recibidas sean correctas para su posterior uso en entrenamiento.

    \item \textbf{Identificación de elementos del entorno}:
    El sistema deberá ser capaz de diferenciar los datos correspondientes al robot y al objetivo. Con esto, se construirán las observaciones con las 
    que se representará el entorno en cada instante.

    \item \textbf{Control de movimiento del robot}:
    El sistema deberá ser capaz de enviar instrucciones al robot Unitree Go2 con las cuales realizará las acciones designadas previamente. Estas instrucciones 
    corresponden con funciones de alto nivel del \Gls{sdk} del robot y permitirán el movimiento de este a través del entorno.

    \item \textbf{Generación y almacenamiento de datos de entrenamiento}:
    El sistema deberá permitir almacenar los datos asociados a las experiencias de interactuar con el entorno de manera adecuada para su posterior uso en 
    entrenamiento. Para el modelo de mundo deberá guardar estos datos en una lista de tuplas con la forma [\( Obs_t \), Accion, \( Obs_{t+1} \)], mientras que para el modelo de 
    utilidad deberán ser de la forma [\( Obs_t \), Utilidad] (esto se explicará con más detalle en el apartado de desarrollo, sección \ref{sec:desviaciones_planificacion}, 
    página \pageref{sec:desviaciones_planificacion}).

    \item \textbf{Entrenamiento del modelo de mundo}:
    El sistema deberá ser capaz de entrenar un modelo de mundo que permita predecir el estado siguiente del entorno a partir del estado actual y de la acción a 
    ejecutar por el robot.

    \item \textbf{Entrenamiento del modelo de utilidad}:
    El sistema deberá permitir entrenar un modelo de utilidad con el cual predecir el valor de efectividad que tendrá un determinado estado en función del objetivo 
    que se desea alcanzar.

    \item \textbf{Evaluación de los modelos entrenados}:
    El sistema deberá incluir métodos para validar y evaluar los modelos entrenados, mostrar gráficas y métricas adecuadas para comprobar si las predicciones 
    son suficientemente correctas.

    \item \textbf{Realizar pruebas en entorno real}:
    El sistema deberá ser capaz de ejecutar pruebas reales con el robot con el fin de comprobar su comportamiento en situaciones de uso reales.
    
\end{itemize}


\subsection{Requisitos no funcionales}
\label{sec:req_no_funcionales}

Los requisitos no funcionales describen las propiedades de calidad y las restricciones que debe cumplir el sistema durante su funcionamiento.
Para este proyecto se definieron los siguientes:

\begin{itemize}
    \item \textbf{Precisión en la obtención de la posición}:
    La captura de posición del robot y del objetivo deben ser lo suficientemente precisos tanto para que los datos puedan ser utilizados para crear modelos 
    funcionales como para la validación de estos mismos.

    \item \textbf{Transmisión de datos fiable}:
    La transmisión de datos entre Motive y el ordenador receptor deberá ser fiable, con el fin de minimizar posibles fallos o pérdidas en la conexión. 
    Es especialmente relevante la consistencia temporal con las acciones para obtener correctamente los datos de entrenamiento.

    \item \textbf{Seguridad durante las pruebas en entorno real}: 
    El sistema deberá utilizarse siempre en un entorno seguro durante las pruebas con el robot real. Será necesario establecer los límites de movimiento del 
    robot y despejar el entorno para minimizar el riesgo de colisión o caídas durante estas pruebas.

    \item \textbf{Modularidad}:
    El sistema desarrollado deberá estar organizado de manera modular, es decir, separando las distintas fases en códigos independientes (adquisición de 
    datos, entrenamiento, pruebas, control del robot, etc). Esto se realiza con el objetivo de facilitar el mantenimiento y una posible ampliación o modificación 
    posterior, además de adaptarse a la metodología seleccionada, a la planificación inicial y a posibles desviaciones de esta, facilitando su adaptación 
    y justificación.

    \item \textbf{Escalabilidad}:
    Aunque para este proyecto se utilizará únicamente un escenario concreto, el sistema deberá estar diseñado para ser escalable, de forma que pueda 
    adaptarse a nuevos objetivos, conjuntos de datos y modificaciones en los modelos sin necesidad de rehacer por completo el código.

    \item \textbf{Mantenibilidad}:
    El código desarrollado deberá ser claro y estar bien documentado, facilitando la comprensión de este y modificaciones futuras.

    \item \textbf{Eficiencia computacional}:
    El entrenamiento de los modelos deberá presentar un coste computacional razonable para los recursos disponibles y los experimentos deberán poder completarse en 
    un tiempo asequible. Esto quiere decir, deberá poder ejecutarse con un hardware similar al utilizado por el alumno y en un 
    tiempo (no tengo muy claro cómo poner esto).

    \item \textbf{Trazabilidad}:  
    El sistema deberá permitir guardar y consultar posteriormente los datos obtenidos de las pruebas realizadas, con el fin de comparar y evaluar los 
    distintos modelos ejecutados durante el desarrollo.
    
\end{itemize}